% problem_id: S001-lemma-hypercovering-descent-bounded-abelian
% source_id: S001 (Stacks Project)
% source_file: spaces-simplicial.tex
% statement_locator: spaces-simplicial.tex lines 4167--4176
% proof_locator: spaces-simplicial.tex lines 4178--4348
% env: lemma
% id_note: label 在本来源内唯一
% pairing: tight (verified)
% status: complete (statement + author proof verbatim + reason + locators)
%
\begin{lemma}
\label{lemma-hypercovering-descent-bounded-abelian}
Let $\mathcal{C}$ be a site with equalizers and fibre products.
Let $K$ be a hypercovering. For
$E \in D(\mathcal{C})$ the map
$$
E \longrightarrow Ra_*a^{-1}E
$$
is an isomorphism.
\end{lemma}

\begin{proof}
First, let $\mathcal{I}$ be an injective abelian sheaf on $\mathcal{C}$.
Then the spectral sequence of
Lemma \ref{lemma-augmentation-spectral-sequence}
for the sheaf $a^{-1}\mathcal{I}$ degenerates as
$(a^{-1}\mathcal{I})_p = a_p^{-1}\mathcal{I}$
is injective by Lemma \ref{lemma-localize-injective}.
Thus the complex
$$
a_{0, *}a_0^{-1}\mathcal{I} \to
a_{1, *}a_1^{-1}\mathcal{I} \to
a_{2, *}a_2^{-1}\mathcal{I} \to\ldots
$$
computes $Ra_*a^{-1}\mathcal{I}$. By
Lemma \ref{lemma-hypercovering-cech-complex}
this is equal to the complex
$\SheafHom(s(\mathbf{Z}_{F(K)}^\#), \mathcal{I})$.
Because $K$ is a hypercovering, we see that
$s(\mathbf{Z}_{F(K)}^\#)$ is exact in degrees $> 0$ by
Hypercoverings, Lemma \ref{hypercovering-lemma-acyclic-hypercover-sheaves}
applied to the simplicial presheaf $F(K)$.
Since $\mathcal{I}$ is injective, the functor $\SheafHom(-, \mathcal{I})$
is exact and we conclude that
$\SheafHom(s(\mathbf{Z}_{F(K)}^\#), \mathcal{I})$
is exact in positive degrees. We conclude that
$R^pa_*a^{-1}\mathcal{I} = 0$ for $p > 0$.
On the other hand, we have $\mathcal{I} = a_*a^{-1}\mathcal{I}$
by Lemma \ref{lemma-hypercovering-descent-sheaves}.

\medskip\noindent
Bounded case. Let $E \in D^+(\mathcal{C})$.
Choose a bounded below complex $\mathcal{I}^\bullet$ of injectives
representing $E$. By the result of the first paragraph and
Leray's acyclicity lemma
(Derived Categories, Lemma \ref{derived-lemma-leray-acyclicity})
$Ra_*a^{-1}\mathcal{I}^\bullet$
is computed by the complex
$a_*a^{-1}\mathcal{I}^\bullet = \mathcal{I}^\bullet$
and we conclude the lemma is true in this case.

\medskip\noindent
Unbounded case. We urge the reader to skip this, since the argument
is the same as above, except that we use explicit representation
by double complexes to get around convergence issues.
Let $E \in D(\mathcal{C})$.
To show the map $E \to Ra_*a^{-1}E$ is an isomorphism,
it suffices to show for every object $U$ of $\mathcal{C}$ that
$$
R\Gamma(U, E) = R\Gamma(U, Ra_*a^{-1}E)
$$
We will compute both sides and show the map $E \to Ra_*a^{-1}E$
induces an isomorphism. Choose a K-injective
complex $\mathcal{I}^\bullet$ representing $E$. Choose a quasi-isomorphism
$a^{-1}\mathcal{I}^\bullet \to \mathcal{J}^\bullet$
for some K-injective complex $\mathcal{J}^\bullet$ on
$(\mathcal{C}/K)_{total}$. We have
$$
R\Gamma(U, E) = R\Hom(\mathbf{Z}_U^\#, E)
$$
and
$$
R\Gamma(U, Ra_*a^{-1}E) = R\Hom(\mathbf{Z}_U^\#, Ra_*a^{-1}E) =
R\Hom(a^{-1}\mathbf{Z}_U^\#, a^{-1}E)
$$
By Lemma \ref{lemma-simplicial-resolution-augmentation}
we have a quasi-isomorphism
$$
\Big(\ldots \to
g_{2!}(a_2^{-1}\mathbf{Z}_U^\#) \to
g_{1!}(a_1^{-1}\mathbf{Z}_U^\#) \to
g_{0!}(a_0^{-1}\mathbf{Z}_U^\#)\Big)
\longrightarrow
a^{-1}\mathbf{Z}_U^\#
$$
Hence $R\Hom(a^{-1}\mathbf{Z}_U^\#, a^{-1}E)$ is equal to
$$
R\Gamma((\mathcal{C}/K)_{total},
R\SheafHom(
\ldots \to
g_{2!}(a_2^{-1}\mathbf{Z}_U^\#) \to
g_{1!}(a_1^{-1}\mathbf{Z}_U^\#) \to
g_{0!}(a_0^{-1}\mathbf{Z}_U^\#),
\mathcal{J}^\bullet))
$$
By the construction in Cohomology on Sites, Section
\ref{sites-cohomology-section-internal-hom}
and since $\mathcal{J}^\bullet$ is K-injective, we see that
this is represented by the complex of abelian groups with terms
$$
\prod\nolimits_{p + q = n}
\Hom(g_{p!}(a_p^{-1}\mathbf{Z}_U^\#), \mathcal{J}^q) =
\prod\nolimits_{p + q = n}
\Hom(a_p^{-1}\mathbf{Z}_U^\#, g_p^{-1}\mathcal{J}^q)
$$
See Cohomology on Sites, Lemmas
\ref{sites-cohomology-lemma-RHom-into-K-injective} and
\ref{sites-cohomology-lemma-section-RHom-over-U} for more information.
Thus we find that $R\Gamma(U, Ra_*a^{-1}E)$ is computed by
the product total complex $\text{Tot}_\pi(B^{\bullet, \bullet})$
with $B^{p, q} = \Hom(a_p^{-1}\mathbf{Z}_U^\#, g_p^{-1}\mathcal{J}^q)$.
For the other side we argue similarly. First we note that
$$
s(\mathbf{Z}_{F(K)}^\#) \longrightarrow \mathbf{Z}
$$
is a quasi-isomorphism of complexes on $\mathcal{C}$
by Hypercoverings, Lemma \ref{hypercovering-lemma-acyclic-hypercover-sheaves}.
Since $\mathbf{Z}_U^\#$ is a flat sheaf of $\mathbf{Z}$-modules
we see that
$$
s(\mathbf{Z}_{F(K)}^\#) \otimes_\mathbf{Z} \mathbf{Z}_U^\#
\longrightarrow
\mathbf{Z}_U^\#
$$
is a quasi-isomorphism. Therefore
$R\Hom(\mathbf{Z}_U^\#, E)$ is equal to
$$
R\Gamma(\mathcal{C}, R\SheafHom(
s(\mathbf{Z}_{F(K)}^\#) \otimes_\mathbf{Z} \mathbf{Z}_U^\#,
\mathcal{I}^\bullet))
$$
By the construction of $R\SheafHom$ and since $\mathcal{I}^\bullet$
is K-injective, this is represented by the complex of abelian groups
with terms
$$
\prod\nolimits_{p + q = n}
\Hom(\mathbf{Z}^\#_{K_p} \otimes_\mathbf{Z} \mathbf{Z}_U^\#, \mathcal{I}^q)
=
\prod\nolimits_{p + q = n}
\Hom(a_p^{-1}\mathbf{Z}_U^\#, a_p^{-1}\mathcal{I}^q)
$$
The equality of terms follows from the fact that
$\mathbf{Z}^\#_{K_p} \otimes_\mathbf{Z} \mathbf{Z}_U^\# =
a_{p!}a_p^{-1}\mathbf{Z}_U^\#$ by Modules on Sites,
Remark \ref{sites-modules-remark-j-shriek-tensor}.
Thus we find that $R\Gamma(U, E)$ is computed by
the product total complex $\text{Tot}_\pi(A^{\bullet, \bullet})$
with $A^{p, q} = \Hom(a_p^{-1}\mathbf{Z}_U^\#, a_p^{-1}\mathcal{I}^q)$.

\medskip\noindent
Since $\mathcal{I}^\bullet$ is K-injective we see that
$a_p^{-1}\mathcal{I}^\bullet$ is K-injective, see
Lemma \ref{lemma-localize-injective}.
Since $\mathcal{J}^\bullet$ is K-injective we see that
$g_p^{-1}\mathcal{J}^\bullet$ is K-injective, see
Lemma \ref{lemma-restriction-injective-to-component-site}.
Both represent the object $a_p^{-1}E$.
Hence for every $p \geq 0$ the map of complexes
$$
A^{p, \bullet} = \Hom(a_p^{-1}\mathbf{Z}_U^\#, a_p^{-1}\mathcal{I}^\bullet)
\longrightarrow
\Hom(a_p^{-1}\mathbf{Z}_U^\#, g_p^{-1}\mathcal{J}^\bullet) = B^{p, \bullet}
$$
induced by $g_p^{-1}$ applied to the given map
$a^{-1}\mathcal{I}^\bullet \to \mathcal{J}^\bullet$
is a quasi-isomorphisms as these complexes both compute
$$
R\Hom(a_p^{-1}\mathbf{Z}_U^\#, a_p^{-1}E)
$$
By More on Algebra, Lemma \ref{more-algebra-lemma-prod-qis-gives-qis}
we conclude that the right vertical arrow in the commutative diagram
$$
\xymatrix{
R\Gamma(U, E) \ar[r] \ar[d] &
\text{Tot}_\pi(A^{\bullet, \bullet}) \ar[d] \\
R\Gamma(U, Ra_*a^{-1}E) \ar[r] &
\text{Tot}_\pi(B^{\bullet, \bullet})
}
$$
is a quasi-isomorphism. Since we saw above that the horizontal arrows
are quasi-isomorphisms, so is the left vertical arrow.
\end{proof}
